% =====================================================================
%  BUỔI 11 — KẾ HOẠCH GIẢNG DẠY (từ giao-vien/01-ke-hoach-chua-de.md)
%  Sinh tự động bằng scripts/md2tex-chung.py — không sửa tay.
% =====================================================================
\documentclass[11pt, a4paper]{article}
% =====================================================================
%  CAE SHANGHAI — GÓI LỆNH DÙNG CHUNG CHO TÀI LIỆU LUYỆN THI CSCA TOÁN
%  Biên dịch bằng XeLaTeX. Mọi file .tex trong buổi học \input file này.
% =====================================================================

% ---------- NGÔN NGỮ & PHÔNG CHỮ ----------
\usepackage{fontspec}
\usepackage{polyglossia}
\setmainlanguage{vietnamese}
\setotherlanguage{english}

% Chữ Latin (Anh + Việt): TeX Gyre Termes — bản clone Times New Roman,
% đi kèm TeX Live, hỗ trợ đầy đủ dấu tiếng Việt.
\setmainfont{TeX Gyre Termes}
\setsansfont{TeX Gyre Heros}

% Chữ Hán: ưu tiên Source Han Serif SC (思源宋体) — chuẩn học thuật Trung Quốc.
% Nếu máy chưa cài, tự động dùng Songti SC (宋体, có sẵn trên macOS).
% Cả hai đều thuộc họ 宋体 nên diện mạo tài liệu không đổi.
\IfFontExistsTF{Source Han Serif SC}{
  \newcommand{\cjkfamilyname}{Source Han Serif SC}
}{
  \IfFontExistsTF{Songti SC}{
    \newcommand{\cjkfamilyname}{Songti SC}
  }{
    \newcommand{\cjkfamilyname}{Noto Serif CJK SC}
  }
}

% xeCJK tự chuyển font khi gặp ký tự Hán, không cần bọc \cjk{} thủ công.
% Đây là điều kiện bắt buộc: font Latin (TeX Gyre Termes) không có glyph chữ Hán,
% nếu không khai báo xeCJK thì mọi chữ Hán viết trần sẽ bị nuốt khỏi bản in.
\usepackage{xeCJK}
\setCJKfamilyfont{caehan}{\cjkfamilyname}[Scale=0.95]
\setCJKmainfont{\cjkfamilyname}[Scale=0.95]
\setCJKsansfont{\cjkfamilyname}[Scale=0.95]
\setCJKmonofont{\cjkfamilyname}[Scale=0.95]

% \cjk{} ép dùng font Hán trong ngữ cảnh cần thiết (ví dụ trong tiêu đề in đậm
% hoặc trong bảng có định dạng riêng). Dùng \CJKfamily chứ không phải tên font,
% vì tên font chỉ là văn bản — sẽ in ra chữ thay vì đổi font.
\newcommand{\cjk}[1]{{\CJKfamily{caehan}#1}}

% ---------- TOÁN HỌC ----------
\usepackage{amsmath, amssymb, amsthm, mathtools}
\usepackage{unicode-math}
\setmathfont{TeX Gyre Termes Math}

% ---------- BỐ CỤC TRANG ----------
% Lề phải chừa vùng an toàn của letterhead:
%   banner trên chiếm 0–19mm  → top tối thiểu 2,2cm
%   banner dưới chiếm 281–297mm → bottom tối thiểu 2,2cm
% Lề trên 2,6cm (26mm) để tiêu đề đầu trang không đè banner (banner hết ở 19mm).
\usepackage[a4paper, top=2.6cm, bottom=2.4cm, left=2.5cm, right=2.0cm]{geometry}
\usepackage{fancyhdr}
\usepackage{titlesec}
\usepackage{enumitem}
\usepackage{graphicx}
\usepackage{booktabs, array, longtable, tabularx}
\usepackage{xcolor}
\usepackage[most]{tcolorbox}
\usepackage[colorlinks=true, linkcolor=cae_blue, citecolor=cae_blue, urlcolor=cae_blue]{hyperref}

% ---------- BẢNG MÀU THƯƠNG HIỆU ----------
\definecolor{cae_blue}{RGB}{0, 51, 102}
\definecolor{cae_gold}{RGB}{196, 149, 60}
\definecolor{box_bg}{RGB}{245, 247, 250}
\definecolor{box_line}{RGB}{0, 51, 102}

% ---------- TIÊU ĐỀ MỤC ----------
% Đệm đầu trang: banner letterhead chiếm 0–19mm, mà lề trên là 22mm.
% \section* đầu trang (sau \clearpage) sát lề nên phải chừa thêm, nếu không
% sẽ đè lên banner.
\titleformat{\section}
  {\normalfont\large\bfseries\sffamily\color{cae_blue}}{\thesection.}{0.6em}{}
\titleformat{\subsection}
  {\normalfont\normalsize\bfseries\sffamily\color{cae_blue!85!black}}{\thesubsection}{0.6em}{}
\titlespacing*{\section}{0pt}{14pt}{6pt}
\titlespacing*{\subsection}{0pt}{10pt}{4pt}

% Tiêu đề không đánh số ở đầu trang: \section* sau \clearpage nằm sát lề trên
% nên đè lên banner letterhead. Dùng \vspace* ngay TRƯỚC \section* không đủ,
% vì titlesec chèn khoảng cách riêng. Cách chắc chắn là đệm bằng \null.
\newcommand{\sectiondautrang}[1]{\section*{#1}}

% Dạng bài dùng \subsection*{} nên cần định dạng riêng, đậm và có gạch chân màu
\titleformat{\subsubsection}
  {\normalfont\normalsize\bfseries\sffamily\color{cae_blue!85!black}}{\thesubsubsection}{0.6em}{}

% ---------- LETTERHEAD (ảnh nền toàn trang) ----------
% Ảnh letterhead của CAE Shanghai: banner xanh trên/dưới kèm logo chìm giữa.
% Tỉ lệ ảnh 1414x2000 = 0.7070, khớp A4 (0.7071) nên chèn thẳng không méo.
% Đặt ở lớp nền dưới cùng để chữ không bị ảnh che.
% Dùng \AddToShipoutPictureBG (không có dấu *) để áp dụng cho MỌI trang.
\usepackage{eso-pic}
\AddToShipoutPictureBG{%
  \AtPageLowerLeft{%
    \includegraphics[width=\paperwidth,height=\paperheight]{images/letterhead.png}%
  }%
}

% ---------- ĐẦU / CHÂN TRANG ----------
% Letterhead đã có banner trên và dưới, nên bỏ dòng chạy để tránh chồng chữ.
% Số trang đặt trong vùng nội dung, sát trên banner dưới, lề phải.
\pagestyle{fancy}
\setlength{\headheight}{0pt}
\fancyhf{}
\renewcommand{\headrulewidth}{0pt}
\renewcommand{\footrulewidth}{0pt}
\fancyfoot[R]{\raisebox{0.9cm}{\small\color{cae_blue}\thepage}}

% ---------- MÔI TRƯỜNG NỘI DUNG ----------
% Trong tcolorbox, LaTeX đặt parindent = parskip = 0, nên các đoạn văn trong
% khung dính liền nhau. Đặt lại khoảng cách đoạn cho mọi khung.
\newcommand{\khoangcachdoan}{\setlength{\parskip}{0.5em}\setlength{\parindent}{0pt}}
\tcbset{before upper={\khoangcachdoan}}

\newtcolorbox{dinhnghia}[1][]{
  colback=box_bg, colframe=box_line, boxrule=0.6pt, arc=1pt,
  fonttitle=\bfseries, title={Định nghĩa #1}, breakable}
\newtcolorbox{dinhly}[1][]{
  colback=box_bg, colframe=cae_gold, boxrule=0.6pt, arc=1pt,
  fonttitle=\bfseries, title={Tính chất #1}, breakable}
\newtcolorbox{vidu}[1][]{
  colback=white, colframe=black!40, boxrule=0.4pt, arc=1pt,
  fonttitle=\bfseries, title={Ví dụ #1}, breakable}
\newtcolorbox{baitap}[1][]{
  colback=white, colframe=cae_blue!50, boxrule=0.4pt, arc=1pt,
  fonttitle=\bfseries, title={Câu #1}, breakable}

% Ô đáp án: viền xanh đậm hơn, dùng trong file đáp án
\newtcolorbox{dapan}[1][]{
  colback=box_bg, colframe=cae_blue, boxrule=0.6pt, arc=1pt,
  fonttitle=\bfseries, title={Câu #1}, breakable}

% ---------- TÊN CHỨNG MINH ----------
\renewcommand{\proofname}{\normalfont\bfseries Chứng minh}

% ---------- MÔN HỌC ----------
% Đổi được ở file .tex trước khi % =====================================================================
%  CAE SHANGHAI — GÓI LỆNH DÙNG CHUNG CHO TÀI LIỆU LUYỆN THI CSCA TOÁN
%  Biên dịch bằng XeLaTeX. Mọi file .tex trong buổi học \input file này.
% =====================================================================

% ---------- NGÔN NGỮ & PHÔNG CHỮ ----------
\usepackage{fontspec}
\usepackage{polyglossia}
\setmainlanguage{vietnamese}
\setotherlanguage{english}

% Chữ Latin (Anh + Việt): TeX Gyre Termes — bản clone Times New Roman,
% đi kèm TeX Live, hỗ trợ đầy đủ dấu tiếng Việt.
\setmainfont{TeX Gyre Termes}
\setsansfont{TeX Gyre Heros}

% Chữ Hán: ưu tiên Source Han Serif SC (思源宋体) — chuẩn học thuật Trung Quốc.
% Nếu máy chưa cài, tự động dùng Songti SC (宋体, có sẵn trên macOS).
% Cả hai đều thuộc họ 宋体 nên diện mạo tài liệu không đổi.
\IfFontExistsTF{Source Han Serif SC}{
  \newcommand{\cjkfamilyname}{Source Han Serif SC}
}{
  \IfFontExistsTF{Songti SC}{
    \newcommand{\cjkfamilyname}{Songti SC}
  }{
    \newcommand{\cjkfamilyname}{Noto Serif CJK SC}
  }
}

% xeCJK tự chuyển font khi gặp ký tự Hán, không cần bọc \cjk{} thủ công.
% Đây là điều kiện bắt buộc: font Latin (TeX Gyre Termes) không có glyph chữ Hán,
% nếu không khai báo xeCJK thì mọi chữ Hán viết trần sẽ bị nuốt khỏi bản in.
\usepackage{xeCJK}
\setCJKfamilyfont{caehan}{\cjkfamilyname}[Scale=0.95]
\setCJKmainfont{\cjkfamilyname}[Scale=0.95]
\setCJKsansfont{\cjkfamilyname}[Scale=0.95]
\setCJKmonofont{\cjkfamilyname}[Scale=0.95]

% \cjk{} ép dùng font Hán trong ngữ cảnh cần thiết (ví dụ trong tiêu đề in đậm
% hoặc trong bảng có định dạng riêng). Dùng \CJKfamily chứ không phải tên font,
% vì tên font chỉ là văn bản — sẽ in ra chữ thay vì đổi font.
\newcommand{\cjk}[1]{{\CJKfamily{caehan}#1}}

% ---------- TOÁN HỌC ----------
\usepackage{amsmath, amssymb, amsthm, mathtools}
\usepackage{unicode-math}
\setmathfont{TeX Gyre Termes Math}

% ---------- BỐ CỤC TRANG ----------
% Lề phải chừa vùng an toàn của letterhead:
%   banner trên chiếm 0–19mm  → top tối thiểu 2,2cm
%   banner dưới chiếm 281–297mm → bottom tối thiểu 2,2cm
% Lề trên 2,6cm (26mm) để tiêu đề đầu trang không đè banner (banner hết ở 19mm).
\usepackage[a4paper, top=2.6cm, bottom=2.4cm, left=2.5cm, right=2.0cm]{geometry}
\usepackage{fancyhdr}
\usepackage{titlesec}
\usepackage{enumitem}
\usepackage{graphicx}
\usepackage{booktabs, array, longtable, tabularx}
\usepackage{xcolor}
\usepackage[most]{tcolorbox}
\usepackage[colorlinks=true, linkcolor=cae_blue, citecolor=cae_blue, urlcolor=cae_blue]{hyperref}

% ---------- BẢNG MÀU THƯƠNG HIỆU ----------
\definecolor{cae_blue}{RGB}{0, 51, 102}
\definecolor{cae_gold}{RGB}{196, 149, 60}
\definecolor{box_bg}{RGB}{245, 247, 250}
\definecolor{box_line}{RGB}{0, 51, 102}

% ---------- TIÊU ĐỀ MỤC ----------
% Đệm đầu trang: banner letterhead chiếm 0–19mm, mà lề trên là 22mm.
% \section* đầu trang (sau \clearpage) sát lề nên phải chừa thêm, nếu không
% sẽ đè lên banner.
\titleformat{\section}
  {\normalfont\large\bfseries\sffamily\color{cae_blue}}{\thesection.}{0.6em}{}
\titleformat{\subsection}
  {\normalfont\normalsize\bfseries\sffamily\color{cae_blue!85!black}}{\thesubsection}{0.6em}{}
\titlespacing*{\section}{0pt}{14pt}{6pt}
\titlespacing*{\subsection}{0pt}{10pt}{4pt}

% Tiêu đề không đánh số ở đầu trang: \section* sau \clearpage nằm sát lề trên
% nên đè lên banner letterhead. Dùng \vspace* ngay TRƯỚC \section* không đủ,
% vì titlesec chèn khoảng cách riêng. Cách chắc chắn là đệm bằng \null.
\newcommand{\sectiondautrang}[1]{\section*{#1}}

% Dạng bài dùng \subsection*{} nên cần định dạng riêng, đậm và có gạch chân màu
\titleformat{\subsubsection}
  {\normalfont\normalsize\bfseries\sffamily\color{cae_blue!85!black}}{\thesubsubsection}{0.6em}{}

% ---------- LETTERHEAD (ảnh nền toàn trang) ----------
% Ảnh letterhead của CAE Shanghai: banner xanh trên/dưới kèm logo chìm giữa.
% Tỉ lệ ảnh 1414x2000 = 0.7070, khớp A4 (0.7071) nên chèn thẳng không méo.
% Đặt ở lớp nền dưới cùng để chữ không bị ảnh che.
% Dùng \AddToShipoutPictureBG (không có dấu *) để áp dụng cho MỌI trang.
\usepackage{eso-pic}
\AddToShipoutPictureBG{%
  \AtPageLowerLeft{%
    \includegraphics[width=\paperwidth,height=\paperheight]{images/letterhead.png}%
  }%
}

% ---------- ĐẦU / CHÂN TRANG ----------
% Letterhead đã có banner trên và dưới, nên bỏ dòng chạy để tránh chồng chữ.
% Số trang đặt trong vùng nội dung, sát trên banner dưới, lề phải.
\pagestyle{fancy}
\setlength{\headheight}{0pt}
\fancyhf{}
\renewcommand{\headrulewidth}{0pt}
\renewcommand{\footrulewidth}{0pt}
\fancyfoot[R]{\raisebox{0.9cm}{\small\color{cae_blue}\thepage}}

% ---------- MÔI TRƯỜNG NỘI DUNG ----------
% Trong tcolorbox, LaTeX đặt parindent = parskip = 0, nên các đoạn văn trong
% khung dính liền nhau. Đặt lại khoảng cách đoạn cho mọi khung.
\newcommand{\khoangcachdoan}{\setlength{\parskip}{0.5em}\setlength{\parindent}{0pt}}
\tcbset{before upper={\khoangcachdoan}}

\newtcolorbox{dinhnghia}[1][]{
  colback=box_bg, colframe=box_line, boxrule=0.6pt, arc=1pt,
  fonttitle=\bfseries, title={Định nghĩa #1}, breakable}
\newtcolorbox{dinhly}[1][]{
  colback=box_bg, colframe=cae_gold, boxrule=0.6pt, arc=1pt,
  fonttitle=\bfseries, title={Tính chất #1}, breakable}
\newtcolorbox{vidu}[1][]{
  colback=white, colframe=black!40, boxrule=0.4pt, arc=1pt,
  fonttitle=\bfseries, title={Ví dụ #1}, breakable}
\newtcolorbox{baitap}[1][]{
  colback=white, colframe=cae_blue!50, boxrule=0.4pt, arc=1pt,
  fonttitle=\bfseries, title={Câu #1}, breakable}

% Ô đáp án: viền xanh đậm hơn, dùng trong file đáp án
\newtcolorbox{dapan}[1][]{
  colback=box_bg, colframe=cae_blue, boxrule=0.6pt, arc=1pt,
  fonttitle=\bfseries, title={Câu #1}, breakable}

% ---------- TÊN CHỨNG MINH ----------
\renewcommand{\proofname}{\normalfont\bfseries Chứng minh}

% ---------- MÔN HỌC ----------
% Đổi được ở file .tex trước khi % =====================================================================
%  CAE SHANGHAI — GÓI LỆNH DÙNG CHUNG CHO TÀI LIỆU LUYỆN THI CSCA TOÁN
%  Biên dịch bằng XeLaTeX. Mọi file .tex trong buổi học \input file này.
% =====================================================================

% ---------- NGÔN NGỮ & PHÔNG CHỮ ----------
\usepackage{fontspec}
\usepackage{polyglossia}
\setmainlanguage{vietnamese}
\setotherlanguage{english}

% Chữ Latin (Anh + Việt): TeX Gyre Termes — bản clone Times New Roman,
% đi kèm TeX Live, hỗ trợ đầy đủ dấu tiếng Việt.
\setmainfont{TeX Gyre Termes}
\setsansfont{TeX Gyre Heros}

% Chữ Hán: ưu tiên Source Han Serif SC (思源宋体) — chuẩn học thuật Trung Quốc.
% Nếu máy chưa cài, tự động dùng Songti SC (宋体, có sẵn trên macOS).
% Cả hai đều thuộc họ 宋体 nên diện mạo tài liệu không đổi.
\IfFontExistsTF{Source Han Serif SC}{
  \newcommand{\cjkfamilyname}{Source Han Serif SC}
}{
  \IfFontExistsTF{Songti SC}{
    \newcommand{\cjkfamilyname}{Songti SC}
  }{
    \newcommand{\cjkfamilyname}{Noto Serif CJK SC}
  }
}

% xeCJK tự chuyển font khi gặp ký tự Hán, không cần bọc \cjk{} thủ công.
% Đây là điều kiện bắt buộc: font Latin (TeX Gyre Termes) không có glyph chữ Hán,
% nếu không khai báo xeCJK thì mọi chữ Hán viết trần sẽ bị nuốt khỏi bản in.
\usepackage{xeCJK}
\setCJKfamilyfont{caehan}{\cjkfamilyname}[Scale=0.95]
\setCJKmainfont{\cjkfamilyname}[Scale=0.95]
\setCJKsansfont{\cjkfamilyname}[Scale=0.95]
\setCJKmonofont{\cjkfamilyname}[Scale=0.95]

% \cjk{} ép dùng font Hán trong ngữ cảnh cần thiết (ví dụ trong tiêu đề in đậm
% hoặc trong bảng có định dạng riêng). Dùng \CJKfamily chứ không phải tên font,
% vì tên font chỉ là văn bản — sẽ in ra chữ thay vì đổi font.
\newcommand{\cjk}[1]{{\CJKfamily{caehan}#1}}

% ---------- TOÁN HỌC ----------
\usepackage{amsmath, amssymb, amsthm, mathtools}
\usepackage{unicode-math}
\setmathfont{TeX Gyre Termes Math}

% ---------- BỐ CỤC TRANG ----------
% Lề phải chừa vùng an toàn của letterhead:
%   banner trên chiếm 0–19mm  → top tối thiểu 2,2cm
%   banner dưới chiếm 281–297mm → bottom tối thiểu 2,2cm
% Lề trên 2,6cm (26mm) để tiêu đề đầu trang không đè banner (banner hết ở 19mm).
\usepackage[a4paper, top=2.6cm, bottom=2.4cm, left=2.5cm, right=2.0cm]{geometry}
\usepackage{fancyhdr}
\usepackage{titlesec}
\usepackage{enumitem}
\usepackage{graphicx}
\usepackage{booktabs, array, longtable, tabularx}
\usepackage{xcolor}
\usepackage[most]{tcolorbox}
\usepackage[colorlinks=true, linkcolor=cae_blue, citecolor=cae_blue, urlcolor=cae_blue]{hyperref}

% ---------- BẢNG MÀU THƯƠNG HIỆU ----------
\definecolor{cae_blue}{RGB}{0, 51, 102}
\definecolor{cae_gold}{RGB}{196, 149, 60}
\definecolor{box_bg}{RGB}{245, 247, 250}
\definecolor{box_line}{RGB}{0, 51, 102}

% ---------- TIÊU ĐỀ MỤC ----------
% Đệm đầu trang: banner letterhead chiếm 0–19mm, mà lề trên là 22mm.
% \section* đầu trang (sau \clearpage) sát lề nên phải chừa thêm, nếu không
% sẽ đè lên banner.
\titleformat{\section}
  {\normalfont\large\bfseries\sffamily\color{cae_blue}}{\thesection.}{0.6em}{}
\titleformat{\subsection}
  {\normalfont\normalsize\bfseries\sffamily\color{cae_blue!85!black}}{\thesubsection}{0.6em}{}
\titlespacing*{\section}{0pt}{14pt}{6pt}
\titlespacing*{\subsection}{0pt}{10pt}{4pt}

% Tiêu đề không đánh số ở đầu trang: \section* sau \clearpage nằm sát lề trên
% nên đè lên banner letterhead. Dùng \vspace* ngay TRƯỚC \section* không đủ,
% vì titlesec chèn khoảng cách riêng. Cách chắc chắn là đệm bằng \null.
\newcommand{\sectiondautrang}[1]{\section*{#1}}

% Dạng bài dùng \subsection*{} nên cần định dạng riêng, đậm và có gạch chân màu
\titleformat{\subsubsection}
  {\normalfont\normalsize\bfseries\sffamily\color{cae_blue!85!black}}{\thesubsubsection}{0.6em}{}

% ---------- LETTERHEAD (ảnh nền toàn trang) ----------
% Ảnh letterhead của CAE Shanghai: banner xanh trên/dưới kèm logo chìm giữa.
% Tỉ lệ ảnh 1414x2000 = 0.7070, khớp A4 (0.7071) nên chèn thẳng không méo.
% Đặt ở lớp nền dưới cùng để chữ không bị ảnh che.
% Dùng \AddToShipoutPictureBG (không có dấu *) để áp dụng cho MỌI trang.
\usepackage{eso-pic}
\AddToShipoutPictureBG{%
  \AtPageLowerLeft{%
    \includegraphics[width=\paperwidth,height=\paperheight]{images/letterhead.png}%
  }%
}

% ---------- ĐẦU / CHÂN TRANG ----------
% Letterhead đã có banner trên và dưới, nên bỏ dòng chạy để tránh chồng chữ.
% Số trang đặt trong vùng nội dung, sát trên banner dưới, lề phải.
\pagestyle{fancy}
\setlength{\headheight}{0pt}
\fancyhf{}
\renewcommand{\headrulewidth}{0pt}
\renewcommand{\footrulewidth}{0pt}
\fancyfoot[R]{\raisebox{0.9cm}{\small\color{cae_blue}\thepage}}

% ---------- MÔI TRƯỜNG NỘI DUNG ----------
% Trong tcolorbox, LaTeX đặt parindent = parskip = 0, nên các đoạn văn trong
% khung dính liền nhau. Đặt lại khoảng cách đoạn cho mọi khung.
\newcommand{\khoangcachdoan}{\setlength{\parskip}{0.5em}\setlength{\parindent}{0pt}}
\tcbset{before upper={\khoangcachdoan}}

\newtcolorbox{dinhnghia}[1][]{
  colback=box_bg, colframe=box_line, boxrule=0.6pt, arc=1pt,
  fonttitle=\bfseries, title={Định nghĩa #1}, breakable}
\newtcolorbox{dinhly}[1][]{
  colback=box_bg, colframe=cae_gold, boxrule=0.6pt, arc=1pt,
  fonttitle=\bfseries, title={Tính chất #1}, breakable}
\newtcolorbox{vidu}[1][]{
  colback=white, colframe=black!40, boxrule=0.4pt, arc=1pt,
  fonttitle=\bfseries, title={Ví dụ #1}, breakable}
\newtcolorbox{baitap}[1][]{
  colback=white, colframe=cae_blue!50, boxrule=0.4pt, arc=1pt,
  fonttitle=\bfseries, title={Câu #1}, breakable}

% Ô đáp án: viền xanh đậm hơn, dùng trong file đáp án
\newtcolorbox{dapan}[1][]{
  colback=box_bg, colframe=cae_blue, boxrule=0.6pt, arc=1pt,
  fonttitle=\bfseries, title={Câu #1}, breakable}

% ---------- TÊN CHỨNG MINH ----------
\renewcommand{\proofname}{\normalfont\bfseries Chứng minh}

% ---------- MÔN HỌC ----------
% Đổi được ở file .tex trước khi \input{cae-style} bằng \renewcommand{\monhoc}{...}
\providecommand{\monhoc}{Môn Toán}

% ---------- TRANG BÌA GỌN ----------
% #1 = loại tài liệu (TÀI LIỆU CHUẨN BỊ, TÀI LIỆU BUỔI HỌC, ...)
% #2 = số buổi
% #3 = tên buổi tiếng Việt
% #4 = tên buổi tiếng Hán
% #5 = module và tỉ trọng
\newcommand{\trangbia}[5]{%
  \begin{center}
    \vspace*{0.6cm}
    {\large\bfseries\color{cae_blue} CAE SHANGHAI}\\[0.15cm]
    {\footnotesize Tài liệu luyện thi CSCA · \monhoc}\\[0.5cm]
    {\color{cae_gold}\rule{0.55\linewidth}{0.8pt}}\\[0.35cm]
    {\small\bfseries\color{cae_blue!80!black} #1}\\[0.35cm]
    {\Large\bfseries\color{cae_blue} BUỔI #2 — #3}\\[0.2cm]
    {\large\color{cae_blue!85!black} \cjk{#4}}\\[0.25cm]
    \ifx\relax#5\relax\else{\small #5}\\[0.35cm]\fi
    {\color{cae_gold}\rule{0.55\linewidth}{0.8pt}}\\[0.3cm]
    {\footnotesize Biên soạn: Hoàng Dương, CAE Shanghai \quad·\quad \today}
  \end{center}
  \vspace{0.3cm}
}

% ---------- TRANG BÌA CHO ĐỀ KIỂM TRA ----------
% Khác \trangbia ở chỗ không có chữ "BUỔI N" — dùng cho đề thi, không phải bài giảng.
% #1 = loại tài liệu, #2 = tên kỳ thi (Việt), #3 = tên kỳ thi (Hán), #4 = thông số
\newcommand{\biade}[4]{%
  \begin{center}
    \vspace*{0.6cm}
    {\large\bfseries\color{cae_blue} CAE SHANGHAI}\\[0.15cm]
    {\footnotesize Tài liệu luyện thi CSCA · \monhoc}\\[0.5cm]
    {\color{cae_gold}\rule{0.55\linewidth}{0.8pt}}\\[0.35cm]
    {\small\bfseries\color{cae_blue!80!black} #1}\\[0.35cm]
    {\Large\bfseries\color{cae_blue} #2}\\[0.2cm]
    {\large\color{cae_blue!85!black} \cjk{#3}}\\[0.25cm]
    {\small #4}\\[0.35cm]
    {\color{cae_gold}\rule{0.55\linewidth}{0.8pt}}\\[0.3cm]
    {\footnotesize Biên soạn: Hoàng Dương, CAE Shanghai \quad·\quad \today}
  \end{center}
  \vspace{0.3cm}
}

% ---------- TIỆN ÍCH ----------
% Nhãn từ khóa in đậm ở đầu đoạn, dùng cho Từ vựng / Từ khóa nhận dạng / Chú ý
% Khoảng cách đoạn do khung tự đặt (xem \khoangcachdoan ở trên),
% ở đây chỉ cần tách đoạn và bỏ thụt đầu dòng.
\newcommand{\nhan}[1]{\par\noindent\textbf{#1.}\ \ignorespaces}

% Bảng gọn, không đánh số
\newcommand{\bang}[1]{%
  \begin{center}#1\end{center}}

% Đánh dấu ô trống cho học sinh điền
\newcommand{\chotrong}{\makebox[3.2cm]{\dotfill}}
 bằng \renewcommand{\monhoc}{...}
\providecommand{\monhoc}{Môn Toán}

% ---------- TRANG BÌA GỌN ----------
% #1 = loại tài liệu (TÀI LIỆU CHUẨN BỊ, TÀI LIỆU BUỔI HỌC, ...)
% #2 = số buổi
% #3 = tên buổi tiếng Việt
% #4 = tên buổi tiếng Hán
% #5 = module và tỉ trọng
\newcommand{\trangbia}[5]{%
  \begin{center}
    \vspace*{0.6cm}
    {\large\bfseries\color{cae_blue} CAE SHANGHAI}\\[0.15cm]
    {\footnotesize Tài liệu luyện thi CSCA · \monhoc}\\[0.5cm]
    {\color{cae_gold}\rule{0.55\linewidth}{0.8pt}}\\[0.35cm]
    {\small\bfseries\color{cae_blue!80!black} #1}\\[0.35cm]
    {\Large\bfseries\color{cae_blue} BUỔI #2 — #3}\\[0.2cm]
    {\large\color{cae_blue!85!black} \cjk{#4}}\\[0.25cm]
    \ifx\relax#5\relax\else{\small #5}\\[0.35cm]\fi
    {\color{cae_gold}\rule{0.55\linewidth}{0.8pt}}\\[0.3cm]
    {\footnotesize Biên soạn: Hoàng Dương, CAE Shanghai \quad·\quad \today}
  \end{center}
  \vspace{0.3cm}
}

% ---------- TRANG BÌA CHO ĐỀ KIỂM TRA ----------
% Khác \trangbia ở chỗ không có chữ "BUỔI N" — dùng cho đề thi, không phải bài giảng.
% #1 = loại tài liệu, #2 = tên kỳ thi (Việt), #3 = tên kỳ thi (Hán), #4 = thông số
\newcommand{\biade}[4]{%
  \begin{center}
    \vspace*{0.6cm}
    {\large\bfseries\color{cae_blue} CAE SHANGHAI}\\[0.15cm]
    {\footnotesize Tài liệu luyện thi CSCA · \monhoc}\\[0.5cm]
    {\color{cae_gold}\rule{0.55\linewidth}{0.8pt}}\\[0.35cm]
    {\small\bfseries\color{cae_blue!80!black} #1}\\[0.35cm]
    {\Large\bfseries\color{cae_blue} #2}\\[0.2cm]
    {\large\color{cae_blue!85!black} \cjk{#3}}\\[0.25cm]
    {\small #4}\\[0.35cm]
    {\color{cae_gold}\rule{0.55\linewidth}{0.8pt}}\\[0.3cm]
    {\footnotesize Biên soạn: Hoàng Dương, CAE Shanghai \quad·\quad \today}
  \end{center}
  \vspace{0.3cm}
}

% ---------- TIỆN ÍCH ----------
% Nhãn từ khóa in đậm ở đầu đoạn, dùng cho Từ vựng / Từ khóa nhận dạng / Chú ý
% Khoảng cách đoạn do khung tự đặt (xem \khoangcachdoan ở trên),
% ở đây chỉ cần tách đoạn và bỏ thụt đầu dòng.
\newcommand{\nhan}[1]{\par\noindent\textbf{#1.}\ \ignorespaces}

% Bảng gọn, không đánh số
\newcommand{\bang}[1]{%
  \begin{center}#1\end{center}}

% Đánh dấu ô trống cho học sinh điền
\newcommand{\chotrong}{\makebox[3.2cm]{\dotfill}}
 bằng \renewcommand{\monhoc}{...}
\providecommand{\monhoc}{Môn Toán}

% ---------- TRANG BÌA GỌN ----------
% #1 = loại tài liệu (TÀI LIỆU CHUẨN BỊ, TÀI LIỆU BUỔI HỌC, ...)
% #2 = số buổi
% #3 = tên buổi tiếng Việt
% #4 = tên buổi tiếng Hán
% #5 = module và tỉ trọng
\newcommand{\trangbia}[5]{%
  \begin{center}
    \vspace*{0.6cm}
    {\large\bfseries\color{cae_blue} CAE SHANGHAI}\\[0.15cm]
    {\footnotesize Tài liệu luyện thi CSCA · \monhoc}\\[0.5cm]
    {\color{cae_gold}\rule{0.55\linewidth}{0.8pt}}\\[0.35cm]
    {\small\bfseries\color{cae_blue!80!black} #1}\\[0.35cm]
    {\Large\bfseries\color{cae_blue} BUỔI #2 — #3}\\[0.2cm]
    {\large\color{cae_blue!85!black} \cjk{#4}}\\[0.25cm]
    \ifx\relax#5\relax\else{\small #5}\\[0.35cm]\fi
    {\color{cae_gold}\rule{0.55\linewidth}{0.8pt}}\\[0.3cm]
    {\footnotesize Biên soạn: Hoàng Dương, CAE Shanghai \quad·\quad \today}
  \end{center}
  \vspace{0.3cm}
}

% ---------- TRANG BÌA CHO ĐỀ KIỂM TRA ----------
% Khác \trangbia ở chỗ không có chữ "BUỔI N" — dùng cho đề thi, không phải bài giảng.
% #1 = loại tài liệu, #2 = tên kỳ thi (Việt), #3 = tên kỳ thi (Hán), #4 = thông số
\newcommand{\biade}[4]{%
  \begin{center}
    \vspace*{0.6cm}
    {\large\bfseries\color{cae_blue} CAE SHANGHAI}\\[0.15cm]
    {\footnotesize Tài liệu luyện thi CSCA · \monhoc}\\[0.5cm]
    {\color{cae_gold}\rule{0.55\linewidth}{0.8pt}}\\[0.35cm]
    {\small\bfseries\color{cae_blue!80!black} #1}\\[0.35cm]
    {\Large\bfseries\color{cae_blue} #2}\\[0.2cm]
    {\large\color{cae_blue!85!black} \cjk{#3}}\\[0.25cm]
    {\small #4}\\[0.35cm]
    {\color{cae_gold}\rule{0.55\linewidth}{0.8pt}}\\[0.3cm]
    {\footnotesize Biên soạn: Hoàng Dương, CAE Shanghai \quad·\quad \today}
  \end{center}
  \vspace{0.3cm}
}

% ---------- TIỆN ÍCH ----------
% Nhãn từ khóa in đậm ở đầu đoạn, dùng cho Từ vựng / Từ khóa nhận dạng / Chú ý
% Khoảng cách đoạn do khung tự đặt (xem \khoangcachdoan ở trên),
% ở đây chỉ cần tách đoạn và bỏ thụt đầu dòng.
\newcommand{\nhan}[1]{\par\noindent\textbf{#1.}\ \ignorespaces}

% Bảng gọn, không đánh số
\newcommand{\bang}[1]{%
  \begin{center}#1\end{center}}

% Đánh dấu ô trống cho học sinh điền
\newcommand{\chotrong}{\makebox[3.2cm]{\dotfill}}


\begin{document}

\trangbia{KẾ HOẠCH GIẢNG DẠY}{11}{Hình học phẳng, đường thẳng và đường tròn, conic}{几何综合}{}

\section{PHÂN BỔ THỜI GIAN}

\bang{
\begin{tabularx}{\linewidth}{@{\hspace{6pt}}Xcc@{\hspace{6pt}}}
\toprule
\textbf{Hoạt động} & \textbf{Tỉ trọng} & \textbf{Thời lượng} \\
\midrule
Tổng quan và phân loại lỗi nhóm hình học & 10\% & 8 phút \\
Chữa nhanh các câu đúng & 25\% & 20 phút \\
Chữa kỹ các câu sai & 45\% & 36 phút \\
Hệ thống dạng lặp và kỹ thuật nhanh & 20\% & 16 phút \\
\bottomrule
\end{tabularx}}

Học sinh làm đề trước buổi, mang theo phiếu trả lời. Giáo viên thu phiếu, thống kê tỉ lệ đúng từng câu trước khi vào lớp. Bốn khối dưới đây giả định học sinh đã có kết quả.

\section{PHÂN BỐ CÂU THEO NHÓM}

\bang{
\begin{tabularx}{\linewidth}{@{\hspace{6pt}}Xlcc@{\hspace{6pt}}}
\toprule
\textbf{Nhóm} & \textbf{Câu} & \textbf{Số câu} & \textbf{Tỉ trọng đề} \\
\midrule
Hình học phẳng và không gian & 1–15 & 15 & 31\% \\
Đường thẳng và đường tròn & 16–21 & 6 & 13\% \\
Conic & 22–26 & 5 & 10\% \\
Ôn tập tích lũy & 27–48 & 22 & 46\% \\
\bottomrule
\end{tabularx}}

Ba nhóm hình học cộng lại 26 câu, chiếm 54\% đề. Phần ôn tập tích lũy phủ bốn nhóm đã học: tập hợp, bất đẳng thức, hàm số và đạo hàm, lượng giác, dãy số.

\section{KHỐI 1 — TỔNG QUAN VÀ PHÂN LOẠI LỖI NHÓM HÌNH HỌC (8 phút)}

\subsection{Mục tiêu}

Học sinh thấy được nhóm hình học chiếm hơn nửa đề và biết mình mất điểm ở tầng nào: đọc sai đề, nhầm công thức, hay dựng hình sai.

\subsection{Kịch bản}

\textbf{Phút 1–3 · Chiếu bảng phân bố.} Nói rõ 26 câu hình học trong 48 câu. Học sinh thường coi hình học là phần khó nên bỏ qua, nhưng đây là nhóm chiếm tỉ trọng lớn nhất trong đề này.

\textbf{Phút 4–6 · Ba tầng lỗi.} Chiếu bảng lỗi bên dưới, yêu cầu học sinh tự đối chiếu phiếu trả lời và tự xếp lỗi của mình vào một tầng. Không đọc đáp án ở bước này.

\textbf{Phút 7–8 · Chốt ba câu sai nhiều nhất.} Nêu số câu đúng của lớp ở Câu 8, Câu 10, Câu 48. Đây là ba câu vào khối 3.

\subsection{Bảng lỗi hay mắc — nhóm đọc sai đề và nhầm công thức}

\bang{
\begin{tabularx}{\linewidth}{@{\hspace{6pt}}lllX@{\hspace{6pt}}}
\toprule
\textbf{Tầng lỗi} & \textbf{Biểu hiện trong đề này} & \textbf{Câu} & \textbf{Cách cảnh báo} \\
\midrule
Đọc sai đề & Đổi vai trò $x$ và $y$ khi tìm giao điểm với trục & 16, 17 & Đọc to "giao với trục nào" rồi mới thay biến \\
Đọc sai đề & Nhầm tiêu điểm trên $Ox$ với trên $Oy$ & 23, 25 & Mẫu số lớn nhất nằm dưới biến nào thì trục lớn theo biến đó \\
Nhầm công thức & Lấy hệ số $2p$ làm $p$ & 22, 24 & Khoanh tròn số $2p$ trong đề trước khi tính \\
Nhầm công thức & Quên bình phương bán kính khi viết phương trình đường tròn & 20 & Viết $r$ rồi viết $r^2$ thành hai dòng riêng \\
Nhầm công thức & Quên chuẩn hóa mẫu số ở công thức khoảng cách & 2, 4 & Viết mẫu số $\sqrt{A^2+B^2+C^2}$ trước khi chia \\
\bottomrule
\end{tabularx}}

\subsection{Bảng lỗi hay mắc — nhóm dựng hình và bỏ sót trường hợp}

\bang{
\begin{tabularx}{\linewidth}{@{\hspace{6pt}}lXll@{\hspace{6pt}}}
\toprule
\textbf{Tầng lỗi} & \textbf{Biểu hiện trong đề này} & \textbf{Câu} & \textbf{Cách cảnh báo} \\
\midrule
Dựng hình sai & Nhầm trung tuyến với trung đoạn trong bài tọa độ không gian & 8, 9, 10 & Vẽ ba điểm, đánh dấu trung điểm trước khi tính \\
Dựng hình sai & Không nhận ra trung trực dây chung là đường nối hai tâm & 21 & Vẽ hai đường tròn cắt nhau, kẻ đường nối tâm \\
Bỏ sót trường hợp & Chỉ xét phương trình bậc hai ở bài tham số & 42 & Thấy tham số ở hệ số của $x^2$ thì hỏi ngay: nếu nó bằng 0 thì sao \\
\bottomrule
\end{tabularx}}

\section{KHỐI 2 — CHỮA NHANH CÁC CÂU ĐÚNG (20 phút)}

\subsection{Mục tiêu}

Chốt đáp án và cách làm gọn cho các câu lớp làm tốt, đồng thời rút ra công thức dùng lại. Không giảng lại lý thuyết.

\subsection{Nguyên tắc}

Mỗi câu tối đa 40 giây. Giáo viên đọc đáp án, nói một câu về cách nhận dạng, rồi đi tiếp. Chỉ dừng lâu khi có học sinh giơ tay hỏi.

\subsection{Danh sách chữa nhanh}

\bang{
\begin{tabularx}{\linewidth}{@{\hspace{6pt}}lcX@{\hspace{6pt}}}
\toprule
\textbf{Câu} & \textbf{Đáp án} & \textbf{Một câu chốt} \\
\midrule
1 & C & Tâm đọc ngược dấu trong ngoặc, bán kính là căn của vế phải \\
3 & D & Pháp tuyến tỉ lệ thì song song hoặc trùng; so hằng số để phân biệt \\
5 & A & Mặt cầu đường kính $MN$: tâm là trung điểm, bán kính bằng nửa $MN$ \\
6 & C & Đối xứng qua trục $x$: giữ $x$, đổi dấu $y$ và $z$ \\
7 & B & Cao độ bằng $0$ thì điểm thuộc mặt $xOy$ \\
11 & B & Thiết diện mặt cầu luôn là hình tròn \\
12 & D & Pháp tuyến vuông góc với mọi vectơ trong mặt phẳng \\
13 & C & Đối xứng qua trục $y$: giữ $y$, đổi dấu $x$ và $z$ \\
14 & B & $\overrightarrow{OM} = \overrightarrow{AB}$ thì $B = A + \overrightarrow{OM}$ \\
15 & C & Ba cạnh tại một đỉnh đôi một vuông góc buộc đáy vuông và cạnh bên vuông đáy \\
18 & A & Khoảng cách hai điểm: cộng bình phương hiệu từng tọa độ \\
19 & A & Tâm là nửa hệ số đổi dấu, bán kính bình phương bằng tâm bình phương trừ hệ số tự do \\
26 & B & Cùng tiêu điểm thì $c$ không đổi, dùng $a^2 = b^2 + c^2$ \\
27 & C & $A \subseteq B$ thì mọi phần tử của $A$ phải nằm trong $B$ \\
28 & D & $c$ là phần tử của $A$ nên $c \in A$ đúng \\
30 & A & Hiệu hai số hạng chia cho hiệu chỉ số ra công sai \\
31 & D & Hai tập giao nhau: trừ phần đếm hai lần \\
32 & D & Chia hai vế cho số dương giữ chiều bất đẳng thức \\
33 & A & Tiếp tuyến: tính $f(1)$ và $f'(1)$ rồi viết phương trình \\
35 & A & Giải từng tập rồi giao với tập vũ trụ trước khi hợp \\
37 & A & Mệnh đề chứa phủ định: lấy phần bù rồi mới hợp \\
38 & D & Hàm cho theo nhánh: xét riêng từng nhánh rồi hợp nghiệm \\
39 & D & Âm nhân dương ra âm \\
40 & D & Logarit của tích bằng tổng các logarit \\
41 & A & Tìm nguyên hàm rồi thay điều kiện để ra hằng số \\
43 & B & Hiệu hai số hạng chia cho hiệu chỉ số ra công sai \\
44 & A & Xác suất có điều kiện nhân với xác suất nền \\
45 & D & Hàm cos có tập xác định là toàn bộ trục số \\
46 & C & Số hạng thứ $n$ của cấp số nhân: $a_1 q^{n-1}$ \\
47 & B & Bất đẳng thức tam giác chỉ đúng một chiều \\
\bottomrule
\end{tabularx}}

\subsection{Câu cần dừng lâu hơn}

\textbf{Câu 35} hay sai vì học sinh tính $A$ rồi hợp luôn mà quên giao với $U$. Nhắc: tập vũ trụ giới hạn $5 \le x \le 15$, phải cắt trước.

\textbf{Câu 37} là bẫy hai tầng: giải bất phương trình bậc hai, lấy phần bù, rồi hợp. Yêu cầu học sinh viết ra giấy từng bước thay vì nhẩm.

\textbf{Câu 45} dễ chọn nhầm vì học sinh đọc "sai" thành "đúng". Khoanh tròn chữ 错误 trong đề.

\section{KHỐI 3 — CHỮA KỸ CÁC CÂU SAI (36 phút)}

\subsection{Mục tiêu}

Học sinh tự tìm ra chỗ sai của mình, không chỉ nhận đáp án đúng.

\subsection{Kịch bản cố định cho mỗi câu}

1. \textbf{Gọi học sinh nêu cách làm} (1 phút). Không nói đúng sai, chỉ ghi các bước lên bảng.
2. \textbf{Chỉ ra bước sai} (1 phút). Khoanh đúng dòng sai trên bảng.
3. \textbf{Chữa chuẩn} (1,5 phút). Trình bày cách làm gọn.
4. \textbf{Câu tương tự} (1 phút). Cho một biến thể nhỏ, học sinh làm tại chỗ.

\subsection{Danh sách câu chữa kỹ — nhóm ưu tiên cao}

\bang{
\begin{tabularx}{\linewidth}{@{\hspace{6pt}}clclX@{\hspace{6pt}}}
\toprule
\textbf{Ưu tiên} & \textbf{Câu} & \textbf{Đáp án} & \textbf{Lỗi chính} & \textbf{Ghi chú sư phạm} \\
\midrule
1 & 8 & D & Tính sai trung điểm hoặc quên lấy căn & Vẽ hệ trục, đánh dấu ba điểm. Nhấn: trung tuyến nối đỉnh với trung điểm cạnh đối diện \\
2 & 10 & B & Nhầm vai trò các đỉnh khi tìm trung điểm & Cho đổi thứ tự đỉnh, hỏi trung điểm có đổi không \\
3 & 48 & A & Không dựng được công thức tổng quát của $a_n$ & Chỉ cần thử $n = 1, 2$ để thấy $\lambda$ nguyên ở đâu. Đây là câu khó, chỉ chữa hướng tiếp cận \\
4 & 21 & D & Không nhận ra trung trực dây chung là đường nối tâm & Vẽ hình, chỉ rõ tính chất đối xứng \\
5 & 22 & B & Lấy hệ số $12$ làm $p$ thay vì $2p$ & Viết rõ $2p = 12$ thành một dòng riêng \\
6 & 24 & C & Không dùng định nghĩa parabol & Nhắc: khoảng cách đến tiêu điểm bằng khoảng cách đến đường chuẩn \\
7 & 25 & A & Không dùng được hệ thức $c^2 = a^2 + b^2$ & Viết $c$ rồi mới tính tâm sai $e = \dfrac{c}{a}$ \\
8 & 23 & D & Quên dấu trừ ở đường chuẩn & Đường chuẩn luôn nằm về phía đối diện với hướng mở \\
9 & 4 & C & Quên lấy căn ở mẫu số & Viết mẫu số riêng một dòng \\
10 & 16 & C & Nhầm giao với trục $x$ và trục $y$ & Cho $x = 0$ hay $y = 0$ tùy trục cần tìm \\
\bottomrule
\end{tabularx}}

\subsection{Danh sách câu chữa kỹ — nhóm còn lại}

\bang{
\begin{tabularx}{\linewidth}{@{\hspace{6pt}}clclX@{\hspace{6pt}}}
\toprule
\textbf{Ưu tiên} & \textbf{Câu} & \textbf{Đáp án} & \textbf{Lỗi chính} & \textbf{Ghi chú sư phạm} \\
\midrule
11 & 17 & B & Nhầm chiều dấu khi thay tọa độ điểm & Thay trực tiếp vào $y - y_0 = k(x - x_0)$ \\
12 & 20 & C & Quên bình phương bán kính & Viết $r = 3$ rồi $r^2 = 9$ thành hai dòng \\
13 & 2 & A & Rút gọn căn chưa triệt để & Nhân cả tử và mẫu với $\sqrt{6}$ \\
14 & 42 & B & Chỉ xét biệt thức, quên trường hợp bậc nhất & Hỏi: hệ số của $x^2$ có thể bằng $0$ không \\
15 & 36 & C & Không nghĩ ra hướng biến đổi & Chỉ cách nhóm thành tổng ba bình phương \\
16 & 29 & B & Nhầm dấu hàm lượng giác ở góc phần tư thứ ba & Vẽ vòng tròn lượng giác, chỉ góc phần tư thứ ba \\
17 & 34 & A & Lẫn điều kiện biệt thức âm và không âm & Viết rõ $\Delta < 0$ cho trường hợp dương với mọi $x$ \\
18 & 9 & B & Tính sai tọa độ trung điểm & Cùng họ với Câu 8, chữa liền nhau \\
\bottomrule
\end{tabularx}}

\subsection{Thứ tự ưu tiên khi hết giờ}

Nếu chỉ còn 15 phút, chữa bốn câu: \textbf{8, 10, 21, 22}. Đây là bốn câu mà lỗi sai nằm ở nhận dạng chứ không ở tính toán, nên sửa được sẽ kéo theo nhiều câu khác.

Câu 48 là câu khó nhất đề. Nếu lớp yếu, chỉ nêu hướng tiếp cận và cách loại trừ phương án, không chữa hết.

\section{KHỐI 4 — HỆ THỐNG DẠNG LẶP VÀ KỸ THUẬT NHANH (16 phút)}

\subsection{Mục tiêu}

Học sinh gom các câu rời thành bốn dạng, mỗi dạng một công thức chốt.

\subsection{Dạng lặp trong đề}

\bang{
\begin{tabularx}{\linewidth}{@{\hspace{6pt}}llX@{\hspace{6pt}}}
\toprule
\textbf{Dạng} & \textbf{Câu} & \textbf{Công thức chốt} \\
\midrule
Khoảng cách trong không gian & 2, 4 & $d = \dfrac{\lvert Ax_0+By_0+Cz_0+D \rvert}{\sqrt{A^2+B^2+C^2}}$ \\
Trung điểm và trung tuyến & 5, 8, 9, 10 & Trung điểm là trung bình cộng từng tọa độ \\
Đối xứng qua trục và qua mặt & 6, 7, 13 & Trục nào giữ tọa độ trục đó, hai tọa độ còn lại đổi dấu \\
Mặt cầu & 1, 5 & $(x-a)^2+(y-b)^2+(z-c)^2=R^2$ \\
Đường tròn & 19, 20, 21 & Tâm là nửa hệ số đổi dấu; $r^2 = \dfrac{D^2+E^2}{4} - F$ \\
Parabol & 22, 24 & $x^2 = 2py$ có tiêu điểm $\left(0, \dfrac{p}{2}\right)$ \\
Hypebol & 23 & $c^2 = a^2 + b^2$, trục thực theo biến mang dấu dương \\
Ellipse & 25, 26 & $c^2 = \lvert a^2 - b^2 \rvert$, trục lớn theo mẫu số lớn hơn \\
\bottomrule
\end{tabularx}}

\subsection{Kỹ thuật nhanh cho phòng thi}

\textbf{Kỹ thuật 1 — So mẫu số để tìm trục lớn.} Với ellipse, không cần nhớ công thức: mẫu số lớn hơn nằm dưới biến nào thì trục lớn theo biến đó. Áp dụng cho Câu 25 và Câu 26.

\textbf{Kỹ thuật 2 — Đọc dấu để tìm trục thực.} Với hypebol, dấu dương đứng trước biến nào thì trục thực theo biến đó. Áp dụng cho Câu 23.

\textbf{Kỹ thuật 3 — Hệ số $2p$ không phải $p$.} Với parabol, luôn viết riêng dòng $2p = $ hệ số, giải ra $p$, rồi mới tìm tiêu điểm. Áp dụng cho Câu 22 và Câu 24.

\textbf{Kỹ thuật 4 — Đối xứng theo tên trục.} Trục nào có tên trong câu hỏi thì tọa độ theo trục đó giữ nguyên, các tọa độ còn lại đổi dấu. Áp dụng cho Câu 6, 7, 13.

\textbf{Kỹ thuật 5 — Câu hỏi giao trục.} Muốn tìm giao với trục $y$ thì cho $x = 0$; muốn tìm giao với trục $x$ thì cho $y = 0$. Áp dụng cho Câu 16.

\subsection{Chốt buổi (3 phút cuối, nằm trong khối 4)}

Viết ba dòng lên góc bảng và để nguyên khi học sinh ra về:

1. \textbf{Mẫu số lớn hơn quyết định trục lớn của ellipse.}
2. \textbf{Dấu dương quyết định trục thực của hypebol.}
3. \textbf{Thấy tham số ở hệ số của $x^2$, hỏi ngay: nếu nó bằng 0 thì sao?}

Dặn dò: buổi 12 học về vectơ và số phức. Học sinh giữ lại phiếu trả lời buổi này để đối chiếu khi làm đề tổng hợp.

\section{GHI CHÚ SAU BUỔI HỌC}

\emph{(Điền sau khi dạy xong, dùng để điều chỉnh buổi sau và các buổi lặp lại.)}

\bang{
\begin{tabularx}{\linewidth}{@{\hspace{6pt}}Xl@{\hspace{6pt}}}
\toprule
\textbf{Nội dung} & \textbf{Ghi nhận} \\
\midrule
Thời gian thực tế từng khối &  \\
Câu học sinh sai nhiều nhất &  \\
Câu hỏi học sinh hỏi nhiều &  \\
Dạng cần nhắc lại ở buổi 12 &  \\
Điều chỉnh cho lần dạy sau &  \\
\bottomrule
\end{tabularx}}
\vspace{0.8cm}
\begin{center}
{\footnotesize\color{cae_blue} Tài liệu nội bộ, CAE SHANGHAI}
\end{center}

\end{document}
